\documentclass[11pt]{article}
\usepackage[margin=1in]{geometry}
\usepackage{amsmath,amssymb,amsthm,mathtools}
\usepackage[colorlinks=true,linkcolor=blue,citecolor=blue,urlcolor=blue]{hyperref}

\newtheorem{theorem}{Theorem}
\newtheorem{lemma}{Lemma}
\newtheorem{remark}{Remark}
\newcommand{\R}{\mathbb{R}}
\newcommand{\eR}{\overline{\mathbb{R}}}
\newcommand{\cO}{\mathcal{O}}
\newcommand{\E}{\mathcal{E}}
\newcommand{\D}{\mathcal{D}}

\title{Snug Truncs for Countably Cozero-Presented Compactifications}
\author{Run \texttt{fa\_banach\_001}, lane 7}
\date{August 11, 2026}

\begin{document}
\maketitle

\paragraph{Status.}
This packet gives a candidate substantial partial result toward Ball's
Conjecture 6.1.4.  The proof is complete for the stated spatial,
countably cozero-presented class.  It does not settle arbitrary
compactifications of Lindel\"of frames.

\section{The question and the corrected formulation}

Let \(q:M\to L\) be a compactification: a dense surjective pointed frame
homomorphism with compact regular domain.  Ball defines \(\E_0q\) to be the
pointed real maps on \(L\) that have an almost-finite extended-real lift to
\(M\).  A trunc \(G\subseteq\E_0q\) is \emph{snug} when its bounded cozero and
con sets generate \(M\), and when the image sublocale \(q_*L\) is exactly the
intersection of the open sublocales determined by the finite domains of the
lifts of all \(g\in G\) \cite[Definition preceding Theorem 3.3.2]{Ball2021}.

The published conjecture asks whether every compactification \(q:M\to L\) of
a Lindel\"of frame has such a snug trunc \cite[Conjecture 6.1.4]{Ball2021}.
The arXiv v1 formula accidentally ranges the last intersection over
\(g\in\overline G\), the range of truncation.  The published definition
corrects this to \(g\in G\), in agreement with the proof of the representation
theorem.  We address the corrected published question.

\section{The countable-cozero theorem}

We use the spatial notation most natural for the construction.  If
\((X,*)\) is compact Hausdorff pointed and \(Y\subseteq X\) is a dense
pointed subspace, restriction gives a compactification
\[
 q:\cO_*(X)\longrightarrow\cO_*(Y).
\]
For a continuous \(\eR\)-valued function \(f\), write
\(\operatorname{dom}f=f^{-1}(\R)\).

\begin{theorem}[countable cozero presentation]\label{thm:main}
Let \((X,*)\) be compact Hausdorff pointed, let \(Y\subseteq X\) be dense
with \(*\in Y\), and let \(q:\cO_*(X)\to\cO_*(Y)\) be restriction.  Suppose
there are cozero opens
\[
 X\supseteq U_1\supseteq U_2\supseteq\cdots\supseteq Y
\]
such that, as sublocales of \(\cO_*(X)\),
\[
 q_*\cO_*(Y)=\bigcap_{n\geq1}\bigl(U_n\to\cO_*(X)\bigr).
\]
Then \(q\) admits a snugly embedded trunc \(G\subseteq\E_0q\).
\end{theorem}

The conclusion therefore covers the spatial part of the conjecture whenever
the compactification sublocale has a countable cozero-neighborhood
presentation.  Finite presentations and dense cozero subspaces are included.

\section{A hierarchy of pole functions}

Choose continuous functions \(h_n:X\to[0,1]\) such that
\(U_n=\{h_n>0\}\).  Since \(*\in Y\), every \(h_n(*)>0\).  Define continuous
functions \(R_n:X\to[1,\infty]\) recursively by
\begin{equation}\label{eq:poles}
 R_1=\frac1{h_1},
 \qquad
 R_{n+1}=\frac{R_n^2}{h_{n+1}},
\end{equation}
with the natural value \(\infty\) when a denominator vanishes.  The map
\((r,t)\mapsto r^2/t\), from \([1,\infty]\times[0,1]\) to
\([1,\infty]\), is continuous with this convention.  Induction gives
\[
 \operatorname{dom}R_n=U_n.
\]

The decisive feature is strict asymptotic separation.  From
\[
 \frac{R_n}{R_{n+1}}=\frac{h_{n+1}}{R_n}
\]
we see that \(R_n/R_{n+1}\) extends continuously by zero on
\(X\setminus U_{n+1}\).  Multiplying these ratios yields
\begin{equation}\label{eq:ratio}
 \frac{R_i}{R_j}\ \text{extends continuously by zero on }X\setminus U_j,
 \qquad i<j.
\end{equation}
Thus later poles dominate every earlier pole at their common boundary.

Put
\[
 g_n=(R_n-R_n(*))|_Y
\]
and let
\[
 B=\{b|_Y:b\in C(X,\R),\ b(*)=0\}.
\]
Finally, let \(G\) be the vector sublattice of \(C(Y,\R)\) generated by
\(B\cup\{g_n:n\geq1\}\).

\begin{lemma}[finite asymptotic-scale calculus]\label{lem:scale}
Every element of \(G\) has a continuous lift
\(\widetilde g:X\to\eR\), finite on \(Y\), with
\(\widetilde g(*)=0\).
\end{lemma}

\begin{proof}
First consider a linear form involving only the first \(N\) poles,
\[
 \ell=b+\sum_{i=1}^N a_i g_i,
 \qquad b\in B.
\]
If \(k\) is the largest index for which \(a_k\neq0\), then
\eqref{eq:ratio} gives
\[
 \frac{\ell}{R_k}\longrightarrow a_k
 \quad\text{continuously on }X\setminus U_k.
\]
Hence \(\ell\) extends there with value \(+\infty\) or \(-\infty\)
according to the sign of \(a_k\); if every \(a_i=0\), it already extends as
the bounded function \(b\).

Every finite lattice-linear expression can, by distributivity, be written as
a finite supremum of finite infima of such linear forms.  At a boundary point,
compare two forms by reading their pole-coefficient vectors in the order
\[
 (a_N,a_{N-1},\ldots,a_1).
\]
The largest index at which the vectors differ determines their order near the
point by \eqref{eq:ratio}.  If all coefficients belonging to poles infinite at
that point agree, the common divergent part cancels and only bounded continuous
remainders remain.  Therefore finite maxima and minima have continuous
extended-real limits as well.  This supplies the asserted lift of every
element of \(G\).  All generators vanish at \(*\), so every lift does also.
\end{proof}

\section{Proof of Theorem \ref{thm:main}}

For \(g\in G\), Lemma \ref{lem:scale} gives
\(\widetilde g\in\D_0X\).  Restricting its inverse-image frame map to \(Y\)
gives the element represented by \(g\), so the defining square for
\(\E_0q\) commutes.  Hence \(G\subseteq\E_0q\).

The set \(G\) is already a vector lattice.  It is also closed under the
truncation operation.  Indeed, for \(g\geq0\), compose its lift with the
continuous clamp
\[
 \tau:\eR\longrightarrow[0,1],
 \qquad \tau(t)=\min\{\max\{t,0\},1\}.
\]
The function \(\tau\circ\widetilde g\) is bounded and continuous on \(X\),
and it vanishes at \(*\).  Consequently
\(\overline g=(\tau\circ\widetilde g)|_Y\) belongs to \(B\subseteq G\).
Thus \(G\) is a trunc in \(\E_0q\).

It remains to verify snugness.  First, the cozero and con sets arising from
the bounded functions in \(B\) generate \(\cO(X)\).  If \(x\neq *\) lies in
an open \(V\), compact Hausdorff normality gives a function
\(b:X\to[0,1]\), with \(b(*)=0\), such that
\(x\in\operatorname{coz}b\subseteq V\).  If \(*\in V\), it gives a
function \(b:X\to[0,1]\), with \(b(*)=0\) and \(b=1\) on \(X\setminus V\),
so that \(*\in\{b<1\}\subseteq V\).  These functions are fixed by
truncation, and hence the required \(\operatorname{coz}\overline G\) and
\(\operatorname{con}\overline G\) sets join-generate \(\cO(X)\).

Second, the lift of \(g_n\) is \(R_n-R_n(*)\), so
\[
 \operatorname{dom}\widetilde g_n=U_n.
\]
Every lift of an element of \(G\) is finite on the image sublocale, while the
displayed generators already give
\[
 \bigcap_{g\in G}
 \bigl(\operatorname{dom}\widetilde g\to\cO_*(X)\bigr)
 \subseteq
 \bigcap_{n\geq1}\bigl(U_n\to\cO_*(X)\bigr)
 =q_*\cO_*(Y).
\]
The reverse inclusion follows because every \(g\in G\) is finite on \(Y\).
The finite-domain intersection is therefore exact.  Both snugness conditions
hold.

\section{Why the unrestricted upgrade remains open}

The construction succeeds because every finite expression has a largest pole
index.  A successor pole can always be made to dominate a previous master
scale using \eqref{eq:poles}; countable stages can be arranged in a sequence.
For a compactification whose cozero neighborhoods have uncountable cofinality,
there need not be a continuous real-valued master pole dominating all earlier
directions at a limit stage.  Lindel\"ofness of \(L\) guarantees an intersection
by cozero sublocales, but it does not give a countable cozero presentation in a
fixed compactification.  Arbitrarily adjoining reciprocal poles is invalid
because differences of two incomparable infinities need not extend.

Bounded searches for the exact conjecture, title, author, and Madden/snug
terminology found the source but no later primary-source resolution or explicit
duplicate of Theorem \ref{thm:main}.  The unrestricted frame-theoretic
conjecture should therefore remain marked open pending expert review.

\begin{thebibliography}{9}
\bibitem{Ball2019}
R. N. Ball,
\emph{Structural aspects of truncated archimedean vector lattices: simple
elements, good sequences}, arXiv:1906.00439 (2019),
\url{https://arxiv.org/abs/1906.00439}.

\bibitem{Ball2021}
R. N. Ball,
\emph{Structural aspects of truncated archimedean vector lattices: good
sequences, simple elements}, Comment. Math. Univ. Carolin. 62 (2021), 95--134,
\href{https://doi.org/10.14712/1213-7243.2021.009}{doi:10.14712/1213-7243.2021.009}.
\end{thebibliography}

\end{document}
